\section{PMD-Mean 的精确 OPD 同步动力学}
\label{dyn:mean-dynamics}

本节固定 prefix 与有限词表，teacher 和初始 anchor 均严格为正，每轮精确求解 $p_k^{\mathrm M}$ 后令 $q_{k+1}=p_k^{\mathrm M}$。使用单轮实例中的局部记号 $q=q_k$、$r=\log(t/q)$、$\kappa=\KL(q\Vert t)$、$w=p^\star/q$ 与 $\lambda=\KL(q\Vert p^\star)$。关键 KKT 为 $\log w+\lambda w=\eta(r+\kappa)$。单轮 target 与其乘子性质见第~\ref{dyn:mean-static} 节。

\subsection{Mean hard sync 的完整稳定区间}
\label{dyn:sec:mean-hard-sync-stability}

以下仍固定一个 prefix，沿用 $q=q_k(\cdot\mid s)$、$t=t(\cdot\mid s)$、$r_a=\log(t_a/q_a)$、$\kappa=\KL(q\Vert t)$ 和 $w_a=p_a^\star/q_a$。词表有限，teacher 的每个概率均严格为正，步长 $\eta$ 在 outer 轮次间固定。乘子 $\lambda=\KL(q\Vert p^\star)$ 由当前 anchor 决定。跨度统一记为 $b(q,t)=d_H(q,t)=\max_a\log(q_a/t_a)-\min_a\log(q_a/t_a)$，因此逐轮的 $b_k=d_H(q_k,t)$。精确 Mean 的 KKT 写为
\begin{equation}
 \log w_a+\lambda w_a=\eta(r_a+\kappa),
 \qquad \sum_aq_aw_a=1.
 \label{dyn:eq:mean-hard-sync-kkt}
\end{equation}
这里讨论的是每轮精确求解并更新 $q_{k+1}=p^\star$ 的分布空间动力学。步长的稳定区间和单轮 teacher KL 的下降区间是两个不同性质。

\begin{theorem}[Mean hard sync 的全局稳定区间]
\label{dyn:thm:mean-hard-sync-full-stability}
固定严格正的 teacher 和严格正的初始 anchor。若 $0<\eta\le2$，精确 Mean hard sync 满足
\begin{equation*}
 b(q_{k+1},t)<b(q_k,t)
 \quad\text{只要 }q_k\ne t,
 \qquad q_k\longrightarrow t.
\end{equation*}
当 $0<\eta<2$ 时，定义
\begin{equation}
 C_\eta(b)=\max\left\{1-\frac{\eta}{1+\eta b},\,\eta-1\right\}.
 \label{dyn:eq:mean-hard-sync-contraction-factor}
\end{equation}
则 $0\le C_\eta(b)<1$ 对 $b>0$ 成立，而且
\begin{equation}
 b(q_{k+1},t)\le C_\eta(b(q_k,t))b(q_k,t)
 \le C_\eta(b(q_0,t))b(q_k,t).
 \label{dyn:eq:mean-hard-sync-stability-geometric}
\end{equation}
当 $\eta>2$ 时，teacher 固定点局部不稳定。
\end{theorem}

\begin{proof}
首先考虑 $q\ne t$，此时 $\lambda>0$。固定当前轮次的 $\lambda,\kappa$，把 $w_a$ 写为 $r_a$ 的函数 $w(r)$。对式~\eqref{dyn:eq:mean-hard-sync-kkt} 求导得到
\begin{equation*}
 \frac{d\log w(r)}{dr}=\frac{\eta}{1+\lambda w(r)}.
\end{equation*}
令 $v(r)=\log w(r)-r$，则 $v(r_a)=\log(p_a^\star/t_a)$，并且
\begin{equation}
 v'(r)=-1+\frac{\eta}{1+\lambda w(r)}.
 \label{dyn:eq:mean-hard-sync-derivative}
\end{equation}
当 $0<\eta\le2$ 时，$\lambda w(r)>0$ 保证 $-1<v'(r)<1$。在有限区间 $[r_{\min},r_{\max}]$ 上连续性进一步给出 $\sup_r|v'(r)|<1$。对任意动作对应用均值定理，得到严格跨度收缩。

为了得到显式界，归一化保证 $w_{\max}\ge1$，而 $\kappa=-\E_qr\le-r_{\min}$。最大 reward 处的 KKT 因此给出
\begin{equation*}
 \lambda w_{\max}
 =\eta(r_{\max}+\kappa)-\log w_{\max}
 \le\eta(r_{\max}-r_{\min})=\eta b(q,t).
\end{equation*}
结合 $\lambda w(r)>0$，式~\eqref{dyn:eq:mean-hard-sync-derivative} 的值落在
\begin{equation*}
 -1+\frac{\eta}{1+\eta b(q,t)}\le v'(r)<\eta-1.
\end{equation*}
令 $L=-1+\eta/(1+\eta b)$、$U=\eta-1$。由于 $L\le U$，有 $|v'|\le\max\{-L,U\}=C_\eta(b)$；这个写法同时覆盖导数同号和异号的情况，由此得到单步界。$C_\eta$ 关于 $b$ 非减，且当 $0<\eta<2$、$b>0$ 时严格小于一。逐轮归纳即得几何收敛。

边界 $\eta=2$ 的上述显式因子等于一，需要使用精确更新的连续性。归一化给出
\begin{equation*}
 \min_a\log(q_a/t_a)\le0\le\max_a\log(q_a/t_a).
\end{equation*}
因此 $b(q,t)\le b_0$ 蕴含
\begin{equation*}
 t_ae^{-b_0}\le q_a\le t_ae^{b_0}.
\end{equation*}
所有迭代均落在严格正概率单纯形内的紧集上。精确更新关于 $q$ 连续：归一化方程关于 $\lambda$ 严格单调，在 $q\ne t$ 处其根连续；在 $q=t$ 附近，$\lambda\le\log\E_qe^{\eta(r+\kappa)}\to0$，KKT 给出 $p^\star\to t$。

设单调下降的跨度趋于 $b_\infty>0$。在紧集
\begin{equation*}
 \{q:b_\infty\le b(q,t)\le b_0\}
\end{equation*}
上，函数 $b(q,t)-b(p^\star(q),t)$ 连续且处处严格为正，因此存在统一正下界。这与跨度收敛矛盾，故 $b_\infty=0$，并且 $q_k\to t$。若某轮已经到达 $t$，后续轮次保持 $t$。

最后，在 teacher 附近采用下面定义的中心化坐标，更新的线性化为 $y\mapsto(1-\eta)y$。当 $\eta>2$ 时，$|1-\eta|>1$。存在邻域和常数 $c>1$，使每个非零充分小误差的更新范数至少放大 $c$ 倍，因此固定点局部不稳定。
\end{proof}

该定理扩大了 Mean 的 hard sync 稳定范围。下面的 teacher KL 单步下降证明适用于 $0<\eta\le1$；跨度收缩证明则允许 $1<\eta\le2$。在后一范围内，不能把跨度的严格下降直接换成每轮 KL 严格下降。

\subsection{Teacher KL 的单步改进}

\begin{proposition}[精确 Mean 更新的 teacher KL 下降]
\label{dyn:prop:mean-opd-kl-descent}
单轮精确解满足
\begin{equation}
 \KL(p^\star\Vert t)
 \le\kappa-\left(\eta^{-1}-1\right)\KL(p^\star\Vert q)
       -\frac{\lambda}{2\eta}\chi^2(p^\star\Vert q).
 \label{dyn:eq:mean-opd-kl-descent}
\end{equation}
当 $0<\eta\le1$ 且 $q\ne t$ 时，$\KL(p^\star\Vert t)<\KL(q\Vert t)$。
\end{proposition}

\begin{proof}
固定由 Mean 解确定的 $\lambda$，定义
\begin{equation*}
 J_\lambda(p)=\E_pr-\eta^{-1}\KL(p\Vert q)
                    -\frac{\lambda}{2\eta}\chi^2(p\Vert q).
\end{equation*}
对 $p_a$ 求导并加入归一化乘子 $\xi$，其一阶条件为
\begin{equation*}
 r_a-\frac{U_a+1}{\eta}
 -\frac{\lambda}{\eta}\left(\frac{p_a}{q_a}-1\right)+\xi=0.
\end{equation*}
Mean 的 KKT 为 $U_a^\star+\lambda p_a^\star/q_a=\eta(r_a+\kappa)$。取 $\xi=\kappa+(1-\lambda)/\eta$，可验证 $p^\star$ 满足上述条件。$J_\lambda$ 关于 $p$ 的 Hessian 是对角阵，对角项为 $-1/(\eta p_a)-\lambda/(\eta q_a)<0$；因此 $p^\star$ 是其唯一全局最大点。

比较 $p^\star$ 与可行分布 $q$，
\begin{equation*}
 \E_{p^\star}r-\eta^{-1}\KL(p^\star\Vert q)
 -\frac{\lambda}{2\eta}\chi^2(p^\star\Vert q)
 \ge J_\lambda(q)=-\kappa.
\end{equation*}
再用
\begin{equation*}
 \KL(p^\star\Vert t)=\KL(p^\star\Vert q)-\E_{p^\star}r
\end{equation*}
得到式~\eqref{dyn:eq:mean-opd-kl-descent}。若 $q\ne t$，则 $\lambda>0$。此时 $p^\star\ne q$：否则 KKT 要求 $\eta \widetilde r_a=\lambda$ 对所有动作成立，而  $\widetilde r$ 非恒定。这保证 $\chi^2(p^\star\Vert q)>0$，故在 $0<\eta\le1$ 时严格下降。
\end{proof}

\subsection{Log-ratio 跨度与 outer 收敛}

定义逐 prefix 的 log-ratio 跨度
\begin{equation}
 b(q,t)=\max_a\log\frac{q_a}{t_a}
        -\min_a\log\frac{q_a}{t_a}.
 \label{dyn:eq:mean-opd-span}
\end{equation}
该量不受公共归一化常数影响。对于正概率分布，$b(q,t)=0$ 当且仅当 $q=t$。

\begin{theorem}[精确 Mean outer 的跨度收敛率]
\label{dyn:thm:mean-opd-span-convergence}
若 $0<\eta\le1$，则
\begin{equation}
 b(p^\star,t)
 \le\left(1-\frac{\eta}{1+\eta b(q,t)}\right)b(q,t).
 \label{dyn:eq:mean-opd-span-step}
\end{equation}
固定正 teacher、正初始策略和常数 $0<\eta\le1$，每轮精确 Mean inner 后 hard sync。令 $b_k(x)=b(q_k(\cdot\mid s),t(\cdot\mid s))$，则
\begin{equation}
 b_k(x)\le
 \left(1-\frac{\eta}{1+\eta b_0(x)}\right)^k b_0(x)
 \longrightarrow0.
 \label{dyn:eq:mean-opd-span-geometric}
\end{equation}
因此 $q_k(\cdot\mid s)\to t(\cdot\mid s)$。当 $\eta=1$ 时，单步还满足
\begin{equation}
 b_{k+1}(x)\le\frac{b_k(x)^2}{1+b_k(x)}.
 \label{dyn:eq:mean-opd-span-quadratic}
\end{equation}
\end{theorem}

\begin{proof}
固定本轮 $\lambda,\kappa$，将概率比写成 reward 的连续函数
\begin{equation*}
 w(r)=\exp\{\eta(r+\kappa)-W_0(\lambda e^{\eta(r+\kappa)})\}.
\end{equation*}
动作 $a$ 的概率比为 $p_a^\star/q_a=w(r_a)$。它满足
\begin{equation*}
 \log w(r)+\lambda w(r)=\eta(r+\kappa),
 \qquad
 \frac{d\log w(r)}{dr}=\frac{\eta}{1+\lambda w(r)}>0.
\end{equation*}
所以概率比随 reward 严格递增，且其 log-ratio 斜率受到 Lambert-$W$ 修正压缩。对 teacher 的 log-ratio 函数为
\begin{equation*}
 v(r)=\log w(r)-r,
 \qquad v'(r)=-1+\frac{\eta}{1+\lambda w(r)}.
\end{equation*}
在 $0<\eta\le1$ 时，$v'\le0$。设 $r_{\max}=\max_a r_a$、$r_{\min}=\min_a r_a$，并令 $w_{\max}=w(r_{\max})$。由于 $w$ 递增，
\begin{equation*}
 |v'(r)|\le1-\frac{\eta}{1+\lambda w_{\max}}
 \quad (r_{\min}\le r\le r_{\max}).
\end{equation*}
均值定理因此给出
\begin{equation*}
 b(p^\star,t)=v(r_{\min})-v(r_{\max})
 \le\left(1-\frac{\eta}{1+\lambda w_{\max}}\right)
       (r_{\max}-r_{\min}).
\end{equation*}
由归一化 $\E_qw(r)=1$，可知 $w_{\max}\ge1$。其 KKT 与 $\kappa=-\E_qr\le-r_{\min}$ 给出
\begin{equation*}
 \lambda w_{\max}
 =\eta(r_{\max}+\kappa)-\log w_{\max}
 \le\eta(r_{\max}-r_{\min})=\eta b(q,t).
\end{equation*}
代入即可证明式~\eqref{dyn:eq:mean-opd-span-step}。

逐轮应用单步界，先得到 $b_{k+1}\le b_k\le b_0$。函数 $x\mapsto1-\eta/(1+\eta x)$ 递增，因此每轮收缩因子都不超过 $1-\eta/(1+\eta b_0)<1$，归纳得到式~\eqref{dyn:eq:mean-opd-span-geometric}。概率归一化同时保证
\begin{equation*}
 \min_a\log\frac{q_k(a\mid s)}{t(a\mid s)}
 \le0\le
 \max_a\log\frac{q_k(a\mid s)}{t(a\mid s)}.
\end{equation*}
所以每个动作的 log-ratio 绝对值不超过 $b_k(x)$；跨度趋零推出逐动作概率收敛。最后令 $\eta=1$ 即得到式~\eqref{dyn:eq:mean-opd-span-quadratic}。
\end{proof}

$\eta=1$ 的二次上界保证超线性收敛；第~\ref{dyn:sec:mean-hard-sync-cubic-span} 节进一步利用归一化乘子的二阶性质，得到真实局部三次阶数。一般起点上的一次更新仍是 Lambert-$W$ 分布。不同 prefix 的 $b_0(x)$ 可以不同，因此上述逐 prefix 结论不自动给出全 prefix 的统一收敛率。

\subsection{\texorpdfstring{$\eta=1$}{eta=1} 的三次跨度界}
\label{dyn:sec:mean-hard-sync-cubic-span}

\begin{proposition}[Mean 的三次单步上界]
\label{dyn:prop:mean-hard-sync-cubic-span}
当 $\eta=1$ 时，令 $b=b(q,t)$，则
\begin{equation}
 b(p^\star,t)
 \le\frac{e^b b^2}{8+e^b b^2}\,b.
 \label{dyn:eq:mean-hard-sync-cubic-span}
\end{equation}
该界和式~\eqref{dyn:eq:mean-opd-span-quadratic} 同时成立，可以取两者的较小值。特别地，$b\to0$ 时
\begin{equation*}
 b(p^\star,t)\le\left(\frac18+O(b)\right)b^3.
\end{equation*}
\end{proposition}

\begin{proof}
在 $\eta=1$ 时，$0<d\log w/dr\le1$，故
\begin{equation*}
 \max_a\log w_a-\min_a\log w_a\le b.
\end{equation*}
归一化 $\E_qw=1$ 保证 $w_{\min}\le1\le w_{\max}$，从而 $w_{\max}\le e^b$。乘子界给出 $\lambda\le\kappa$。另一方面，由 $\E_qe^r=\sum_at_a=1$，
\begin{equation*}
 \kappa=-\E_qr
 =\log\E_qe^{r-\E_qr}.
\end{equation*}
为写出 Hoeffding 界的常数，令 $F(u)=\log\E_q\exp\{u(r-\E_qr)\}$。有 $F(0)=F'(0)=0$，而 $F''(u)$ 等于指数倾斜分布下 $r$ 的方差。跨度仍为 $b$，故以区间中点为参照有 $F''(u)\le b^2/4$。两次积分给出
\begin{equation*}
 \kappa=F(1)=\int_0^1(1-u)F''(u)\,du\le b^2/8.
\end{equation*}
因此
\begin{equation*}
 \lambda w_{\max}\le e^b b^2/8.
\end{equation*}
式~\eqref{dyn:eq:mean-hard-sync-derivative} 在 $\eta=1$ 时为
\begin{equation*}
 |v'(r)|=\frac{\lambda w(r)}{1+\lambda w(r)}
 \le\frac{e^b b^2}{8+e^b b^2}.
\end{equation*}
再用均值定理便得到结论。$q=t$ 时两边均为零。
\end{proof}

三次界来自两个不同因素：归一化乘子在 teacher 附近是二阶小量，动作间的 reward 差异是一阶小量。仅用 $\lambda w_{\max}\le b$ 会丢掉乘子的二阶性质，因此得到的二次界不能描述真实局部阶数。

\subsection{统一 teacher 坐标下的精细展开}
\label{dyn:sec:mean-hard-sync-local-expansion}

为了比较不同步长，定义同一个 teacher 中心化误差坐标
\begin{equation}
 y_a=\log(q_a/t_a)-\E_t\log(q/t),
 \qquad \E_ty=0,\qquad V=\E_ty^2.
 \label{dyn:eq:mean-hard-sync-centered-coordinate}
\end{equation}
于是 $q_a=t_ae^{y_a}/\E_te^y$。更新后的 $y_a^+$ 采用完全相同的定义，将 $q$ 换成 $p^\star$。记 $\delta=\|y\|_\infty$。以下大 $O$ 常数只依赖固定 teacher、有限词表及固定步长 $\eta$，与接近 teacher 的具体方向无关。

\begin{theorem}[Mean 更新的三阶局部展开]
\label{dyn:thm:mean-hard-sync-local-expansion}
对固定 $\eta>0$，当 $\delta\to0$ 时，
\begin{align}
 \kappa&=\frac12V+O(\delta^3),\notag\\
 \lambda&=\frac{\eta^2}{2}V+O(\delta^3),\notag\\
 y_a^+&=(1-\eta)y_a+\frac{\eta^3}{2}Vy_a+O(\delta^4).
 \label{dyn:eq:mean-hard-sync-local-expansion}
\end{align}
最后一个余项关于动作一致有界。对于 $\eta=1$，有
\begin{equation}
 V^+=\frac14V^3+O(V^{7/2}),
 \qquad
 \frac{\KL(p^\star\Vert t)}{\KL(q\Vert t)^3}\longrightarrow1
 \quad(q\to t,\ q\ne t).
 \label{dyn:eq:mean-hard-sync-kl-cubic}
\end{equation}
\end{theorem}

\begin{proof}
令 $Z=\E_te^y$。$\E_ty=0$ 和有限维 Taylor 展开给出
\begin{equation*}
 Z=1+\frac12V+O(\delta^3),\qquad
 \log Z=\frac12V+O(\delta^3).
\end{equation*}
又因为 $q_a=t_ae^{y_a}/Z$，
\begin{equation*}
 \E_qy=\frac{\E_t[y e^y]}{Z}=V+O(\delta^3).
\end{equation*}
所以 $\kappa=\E_qy-\log Z=V/2+O(\delta^3)$。中心化 reward 恰好等于
\begin{equation*}
 r_a+\kappa=-y_a+\E_qy.
\end{equation*}
记 $g_a=\eta(-y_a+\E_qy)$，则 $\E_qg=0$ 且 $g_a=O(\delta)$。归一化方程为
\begin{equation*}
 \Phi(\lambda,y)
 =\E_q\exp\{g-W_0(\lambda e^g)\}-1=0.
\end{equation*}
在 $(\lambda,y)=(0,0)$，$\Phi=0$ 且 $\partial_\lambda\Phi=-1$。$W_0$ 在零点解析，隐函数定理因此给出局部解析的 $\lambda(y)$。为确定首项，先计算
\begin{equation*}
 \E_qe^g=1+\frac12\E_qg^2+O(\delta^3)
 =1+\frac{\eta^2}{2}V+O(\delta^3).
\end{equation*}
由于 $W_0(z)=z+O(z^2)$，归一化式进一步写为
\begin{equation*}
 0=\E_qe^g-1-\lambda\E_qe^{2G}+O(\lambda^2).
\end{equation*}
隐函数的连续性及此式先给出 $\lambda=O(\delta^2)$，继而得到
\begin{equation*}
 \lambda=\frac{\eta^2}{2}V+O(\delta^3).
\end{equation*}

下面对 KKT 解做 teacher 中心化。记 $C_tf=f-\E_tf$。因公共归一化常数在中心化后消失，
\begin{equation*}
 y^+=y+C_tG-C_tW_0(\lambda e^g)
 =(1-\eta)y-C_tW_0(\lambda e^g).
\end{equation*}
$\lambda$ 在动作间是公共标量，且 $\lambda=O(\delta^2)$，故
\begin{align*}
 C_tW_0(\lambda e^g)
 &=\lambda C_te^g+O(\delta^4),\\
 C_te^g&=C_tG+O(\delta^2)
 =-\eta y+O(\delta^2).
\end{align*}
因此
\begin{equation*}
 y^+=(1-\eta)y+\eta\lambda y+O(\delta^4)
 =(1-\eta)y+\frac{\eta^3}{2}Vy+O(\delta^4).
\end{equation*}

当 $\eta=1$ 时，$y^+=Vy/2+R$，其中 $\|R\|_\infty=O(\delta^4)$。固定 $t_{\min}=\min_at_a>0$ 保证
\begin{equation*}
 t_{\min}\delta^2\le V\le\delta^2.
\end{equation*}
从而 $\|R\|_\infty=O(V^2)$。对平方取 teacher 均值，得到
\begin{equation*}
 V^+=\frac14V^3+V\E_t[yR]+\E_tR^2
 =\frac14V^3+O(V^{7/2}).
\end{equation*}
同一 KL 展开分别用于旧、新分布：
\begin{equation*}
 \KL(q\Vert t)=\frac12V+O(V^{3/2}),
 \qquad
 \KL(p^\star\Vert t)=\frac12V^++O((V^+)^{3/2}).
\end{equation*}
分子首项为 $V^3/8$，分母首项也为 $V^3/8$，由此得到比值极限。
\end{proof}

沿任意固定非恒定方向 $z$，取 $q_\epsilon\propto t\exp(\epsilon z)$。中心化后 $y=\epsilon(z-\E_tz)$。因此 $\eta=1$ 时
\begin{equation*}
 \frac{b(p^\star,t)}{b(q_\epsilon,t)^3}
 \longrightarrow
 \frac{\Var_t(z)}{2(\max_a z_a-\min_a z_a)^2}.
\end{equation*}
Popoviciu 方差界使该极限不超过 $1/8$，与三次跨度上界的局部常数一致。这个方向依赖的跨度常数和方向无关的 KL 三次比值是同一展开的两个结果。

\subsection{边界步长 \texorpdfstring{$\eta=2$}{eta=2} 的幂率}
\label{dyn:sec:mean-hard-sync-boundary-rate}

\begin{theorem}[Mean 边界步长的渐近 KL 率]
\label{dyn:thm:mean-hard-sync-boundary-rate}
固定有限词表、严格正 teacher、严格正初始 anchor，在 $\eta=2$ 下执行精确 Mean hard sync。若某轮到达 teacher，则后续 KL 恒为零。若没有任何有限轮次到达 teacher，则
\begin{equation}
 kV_k\longrightarrow\frac18,
 \qquad k\KL(q_k\Vert t)\longrightarrow\frac1{16}.
 \label{dyn:eq:mean-hard-sync-boundary-rate}
\end{equation}
\end{theorem}

\begin{proof}
全局稳定定理已经给出 $q_k\to t$，从而 $V_k\to0$。三阶展开在 $\eta=2$ 时成为
\begin{equation*}
 y^+=-y+4Vy+R,\qquad \|R\|_\infty=O(V^2).
\end{equation*}
对平方取 $t$ 均值：
\begin{align*}
 V^+
 &=(-1+4V)^2V+2(-1+4V)\E_t[yR]+\E_tR^2\\
 &=V-8V^2+O(V^{5/2}).
\end{align*}
未有限步到达 teacher 意味着所有 $V_k>0$。当 $k$ 足够大时，$V_{k+1}/V_k\to1$，因此
\begin{equation*}
 \frac1{V_{k+1}}-\frac1{V_k}
 =\frac{V_k-V_{k+1}}{V_kV_{k+1}}
 =8+O(\sqrt{V_k})\longrightarrow8.
\end{equation*}
从某个固定轮次开始累加，并对平均增量取极限，得到 $1/(kV_k)\to8$。最后
\begin{equation*}
 \KL(q_k\Vert t)=\frac12V_k+O(V_k^{3/2})
\end{equation*}
给出 $k\KL(q_k\Vert t)\to1/16$。
\end{proof}

Gibbs 更新在同一坐标中恰好满足 $y^+=(1-\eta)y$。当 $\eta=2$ 时，Gibbs 的 $y^+=-y$ 构成两周期，$V$ 恒定。Mean 的 $4Vy$ 修正提供三次阻尼，因而其边界步长收敛，但仅具有 $1/k$ 的 KL 渐近率。当 $\eta=1$ 时，Gibbs 一步到达 teacher，Mean 则保留三阶误差。两者具有相同 teacher 固定点，但精确 outer 更新及其边界行为不同。
